\chapter{A Concrete Naming Certificate for $\CauchyCutUp$}
\label{ch:concrete-instances-cauchy-cut-namecert}
\origin{human}

Following Bishop and Bridges constructive analysis, the $\CauchyCutUp$ packet records the finite bridge by which a located lower--upper cut presentation is read as a regular Cauchy real name before the terminal $\RealUp$ seal is exposed. It sits after the located-cut surface of \autoref{ch:concrete-instances-locatedcut-namecert}, the dyadic arithmetic carrier of \autoref{ch:concrete-instances-dyadicratcore-namecert}, the finite stream-window surface of \autoref{ch:concrete-instances-streamname-namecert}, the regular rational handoff of \autoref{ch:concrete-instances-regseqrat-namecert}, and the terminal real-name seal of \autoref{ch:concrete-instances-real-namecert}. The packet names a displayed cut-to-stream route; it does not identify Dedekind cuts, quotient regular names, selected limits, trichotomy, countable choice, or ambient $\RealUp$ completeness.

\begin{definition}[Cauchy cut carrier]
\label{def:cauchy-cut-carrier}
A $\CauchyCutUp$ carrier is a finite $\BHist$ packet
$$
K=(L,U,M,S,Q,D,R,H,C,P,N)
$$
with the following displayed rows:
\begin{enumerate}
\item $L$ is the located lower-cut row, recording inhabited rounded lower evidence and its rational observation boundary through \autoref{ch:concrete-instances-locatedcut-namecert}.
\item $U$ is the located upper-cut row, recording inhabited rounded upper evidence and the displayed separation from $L$.
\item $M$ is the dyadic Cauchy modulus ledger: for each displayed precision row it supplies a finite dyadic bracket window whose lower and upper endpoints are read from $L$ and $U$.
\item $S$ is the $\StreamNameUp$ finite-window schedule used to ask for the bracket midpoint at a requested precision.
\item $Q$ is the $\RegSeqRatUp$ regular rational source row extracted from the midpoint stream and the modulus ledger.
\item $D$ is the $\DyadicRatCoreUp$ arithmetic row recording endpoint comparison, midpoint formation, and tolerance shrinking for the displayed bracket.
\item $R$ is the terminal $\RealUp$ seal exposed only after $L,U,M,S,Q,D$ have been displayed.
\item $H$ stores componentwise $\hsame$ transport for the cut rows, modulus ledger, stream schedule, regular source, dyadic row, real seal, replay, provenance, and local naming row.
\item $C$ stores displayed $\Cont$ replay of the route
$$
L,U,M,S,D \longmapsto Q \longmapsto R .
$$
\item $P$ is the $\Pkg$ provenance over $\CauchyCutUp$, $\LocatedCutUp$, $\DyadicRatCoreUp$, $\StreamNameUp$, $\RegSeqRatUp$, $\RealUp$, $\BHist$, $\hsame$, $\Cont$, and $\NameCert$ rows.
\item $N$ is the local $\NameCert_{\CauchyCutUp}$ row.
\end{enumerate}
The classifier compares accepted packets only through $L,U,M,S,Q,D,R,H,C,P,N$. It has no coordinate for host cut equality, Dedekind completeness, quotient stream equality, a selected limiting real, trichotomy, countable choice, or ambient $\RealUp$ completeness.
\end{definition}

\begin{theorem}[Cauchy cut carrier admission]
\label{thm:cauchy-cut-carrier-admission}
For every accepted $\CauchyCutUp$ carrier $K=(L,U,M,S,Q,D,R,H,C,P,N)$, the local carrier surface is exactly the displayed lower located-cut row $L$, upper located-cut row $U$, dyadic modulus ledger $M$, stream-window row $S$, regular rational source row $Q$, dyadic arithmetic row $D$, terminal real seal $R$, componentwise transport $H$, displayed replay $C$, provenance $P$, and local naming row $N$.
\end{theorem}
\begin{proof}
Carrier admission is projection from \autoref{def:cauchy-cut-carrier}. The tuple displays the lower and upper located-cut rows before any stream or regular-rational read can occur.

The remaining rows are structural or downstream rows attached to that same carrier. Since the classifier has no other coordinate, an accepted packet can expose no carrier data outside the listed rows.
\end{proof}

\begin{theorem}[Cauchy cut stream extraction]
\label{thm:cauchy-cut-stream-extraction}
Every stream read exported by an accepted $\CauchyCutUp$ packet factors through the displayed route
$$
L,U,M,D \longmapsto S \longmapsto Q .
$$
Thus the extracted regular Cauchy name is obtained from located lower--upper cut brackets and dyadic midpoint arithmetic, not from a quotient stream, selected representative, host real limit, or Dedekind-cut equality.
\end{theorem}
\begin{proof}
Start with the lower and upper cut rows $L$ and $U$. The modulus ledger $M$ chooses only a finite displayed bracket for the requested precision, and the dyadic row $D$ records the midpoint and tolerance calculation inside $\DyadicRatCoreUp$.

The stream row $S$ schedules exactly those finite bracket reads. The regular rational row $Q$ is downstream of $S$ and therefore inherits the same cut and dyadic dependencies. No quotient stream or selected limit appears in the carrier.
\end{proof}

\begin{theorem}[Cauchy cut modulus stability]
\label{thm:cauchy-cut-modulus-stability}
If an accepted $\CauchyCutUp$ carrier is transported componentwise by its $H$ row, the dyadic Cauchy modulus ledger remains attached to the same lower--upper cut bracket route. Later precision reads may refine the displayed bracket through $M$ and $D$, but they cannot replace it by an independent stream selector or by an ambient completeness principle.
\end{theorem}
\begin{proof}
The transport row $H$ acts componentwise on $L,U,M,S,Q,D,R$ and the local naming row. In particular, a transported packet still exposes the same modulus ledger between the located cut rows and the stream schedule.

Refining a precision read uses $M$ again and then replays the dyadic arithmetic row $D$. Since neither $H$ nor $C$ contains a coordinate for a separate selector or completeness theorem, the modulus route is stable under the displayed transport.
\end{proof}

\begin{theorem}[Cauchy cut classifier transport]
\label{thm:cauchy-cut-classifier-transport}
The classifier for $\CauchyCutUp$ transports acceptance only through $L,U,M,S,Q,D,R,H,C,P,N$. Consequently every consumer that reads the packet through \autoref{ch:concrete-instances-locatedcut-namecert}, \autoref{ch:concrete-instances-dyadicratcore-namecert}, \autoref{ch:concrete-instances-streamname-namecert}, \autoref{ch:concrete-instances-regseqrat-namecert}, and \autoref{ch:concrete-instances-real-namecert} receives the same displayed cut-to-regular-name route.
\end{theorem}
\begin{proof}
Inspect the classifier clause in \autoref{def:cauchy-cut-carrier}. It compares accepted packets only by the listed carrier fields and therefore cannot branch on hidden cut equality, host real equality, or quotient stream data.

The cited sibling chapters name the only public read points for those fields: located cuts for $L,U$, dyadic arithmetic for $M,D$, finite streams for $S$, regular rational names for $Q$, and the real seal for $R$. Transport through $H,C,P,N$ preserves that dependency chain for every consumer.
\end{proof}

\begin{theorem}[Cauchy cut real seal nonescape]
\label{thm:cauchy-cut-real-seal-nonescape}
The $\RealUp$ seal of an accepted $\CauchyCutUp$ carrier is public only as the endpoint of the displayed route
$$
L,U,M,S,D,Q \longmapsto R .
$$
No consumer of the packet can extract a completed real by bypassing the located-cut rows, the dyadic Cauchy modulus ledger, the finite stream window, or the regular rational source row.
\end{theorem}
\begin{proof}
The seal row $R$ occurs after $L,U,M,S,Q,D$ in the carrier. The replay row $C$ records the same order, so a consumer must display the located cut rows, the modulus ledger, the stream schedule, the dyadic arithmetic, and the regular rational source before reading the real seal.

The provenance and naming rows $P,N$ attach that route to the local $\NameCert_{\CauchyCutUp}$ surface. Since the carrier excludes selected limits, quotient regular names, Dedekind completeness, countable choice, trichotomy, and ambient $\RealUp$ completeness, the real seal cannot escape the displayed route.
\end{proof}

\falsifiablePrediction{Every accepted $\CauchyCutUp$ read that exports a $\RealUp$ seal factors through displayed located lower--upper cut rows, a dyadic Cauchy modulus ledger, finite stream windows, a regular rational source, dyadic arithmetic, componentwise $\hsame$ transport, displayed $\Cont$ replay, $\Pkg$ provenance, and local naming.}

\independenceWitness{locatedcut, dyadicratcore, streamname, regseqrat, real: $\CauchyCutUp$ is not bijective to any listed sibling because it binds lower--upper located-cut evidence, a dyadic modulus ledger, stream scheduling, regular-rational extraction, terminal real sealing, transport, replay, provenance, and local NameCert rows into one packet.}

\closureat{CauchyCutUp}{seedStr}

\begin{closurestatus}{\CauchyCutUp}
  \theoryclosure{\seedClosure}
  \scopeclosed{\autoref{def:cauchy-cut-carrier}, \autoref{thm:cauchy-cut-carrier-admission}, \autoref{thm:cauchy-cut-stream-extraction}, \autoref{thm:cauchy-cut-modulus-stability}, \autoref{thm:cauchy-cut-classifier-transport}, and \autoref{thm:cauchy-cut-real-seal-nonescape} seed the located-cut-to-regular-Cauchy-name bridge surface for Real and Completion consumers.}
  \formalstatus{\unformalizedV}
  \bridgestatus{none}
  \notclaimed{This chapter does not close Dedekind-cut completeness, quotient regular-name equality, selected limits, host real equality, trichotomy, countable choice, Lean verification, or bridge checking.}
  \upgradepath{Obligation closure requires checked carrier admission, cut-to-stream extraction, modulus stability, classifier transport, real-seal nonescape, provenance discipline, and TasteGate field faithfulness for BEDC.Derived.CauchyCutUp.}
  \constructivestory{A $\CauchyCutUp$ packet is generated from displayed lower and upper located-cut rows, a dyadic Cauchy modulus ledger, finite stream windows, a regular rational extraction row, dyadic arithmetic, a terminal Real seal, componentwise transport, displayed replay, package provenance, and a local NameCert row.}
\end{closurestatus}
